\documentclass[11pt]{amsart}
\usepackage[margin=1.1in]{geometry}
\usepackage{amsmath,amssymb,amsthm,booktabs,enumitem}
\usepackage[hidelinks]{hyperref}

\newtheorem{theorem}{Theorem}[section]
\newtheorem{proposition}[theorem]{Proposition}
\newtheorem{lemma}[theorem]{Lemma}
\newtheorem{corollary}[theorem]{Corollary}
\theoremstyle{definition}
\newtheorem{definition}[theorem]{Definition}
\newtheorem{remark}[theorem]{Remark}
\newtheorem{example}[theorem]{Example}

\numberwithin{equation}{section}

\newcommand{\cP}{\mathcal{P}}
\newcommand{\cQ}{\mathcal{Q}}
\newcommand{\cB}{\mathcal{B}}
\newcommand{\cU}{\mathcal{U}}
\newcommand{\cT}{\mathcal{T}}
\newcommand{\rev}{\operatorname{rev}}
\newcommand{\Z}{\mathbb{Z}}
\newcommand{\R}{\mathbb{R}}
\newcommand{\N}{\mathbb{N}}
\newcommand{\Sinf}{S_\infty}

\title[Palindromic perfect powers]{Palindromic perfect powers: an improvement of the\\ Cilleruelo--Luca--Shparlinski bound}
\date{October 2026}

\begin{document}

\begin{abstract}
For an integer base $g\ge2$ let $\cP_{g,L}$ be the set of palindromes with
$L$ digits in base $g$. Cilleruelo, Luca and Shparlinski proved that for
every $k\ge2$ the number of perfect $k$-th powers in $\cP_{g,L}$ is
$O\bigl((\#\cP_{g,L})^{1/k}\bigr)$. We show that the exponent $1/k$ can be
replaced by $\gamma/k$ with an explicit $\gamma=\gamma(g,k)<1$, for every
base $g$ and every $k$. For instance, one may take $\gamma=0.9674$ for
$g=2$ and every $k$, and the number of $L$-digit palindromic squares in
base $10$ is $O\bigl((\#\cP_{10,L})^{0.4959}\bigr)$.

The proof combines the tree argument of Cilleruelo, Luca and Shparlinski
with a Fourier analysis of the next blocks of digits from both ends. This
reduces the problem to the supremum norm of the Fourier transform of digit
reversal, summed over one frequency, together with an elementary
equidistribution of archimedean phases across the tree. A power saving for
that norm comes from an identity between consecutive digit factors, an idea
that goes back to Col. We obtain good explicit constants from a certified
computation of suprema over windows of consecutive digits.
\end{abstract}

\maketitle

\section{Introduction}

\subsection{Background}
Let $g\ge2$ be an integer. Every positive integer $n$ has a unique expansion
$n=\sum_{i<L}d_ig^i$ with digits $0\le d_i<g$ and $d_{L-1}\ne0$; we call $L$
the length of $n$. The integer $n$ is a \emph{palindrome} (in base $g$) if
$d_i=d_{L-1-i}$ for all $i$. Let $\cP_{g,L}$ be the set of palindromes of
length $L$. A palindrome is determined by its first $\lceil L/2\rceil$
digits, so
\[
\#\cP_{g,L}=(g-1)g^{\lceil L/2\rceil-1}\asymp_g g^{L/2}.
\]
Palindromes are a sparse set, and their arithmetic properties have
attracted much attention; see for instance
\cite{BHS,Col,CLB,TP,DMRSS} and the references there.

Cilleruelo, Luca and Shparlinski \cite{CLS} studied perfect powers among
palindromes. For $k\ge2$ let
\[
\cQ^k_{g,L}=\{n\in\cP_{g,L}:\ n=x^k\text{ for some }x\in\N\}.
\]
They proved
\begin{equation}\label{eq:CLS}
\#\cQ^k_{g,L}\ \ll_{g,k}\ (\#\cP_{g,L})^{1/k}.
\end{equation}
Their argument is elementary. If $x^k$ is a palindrome, its last $t$ digits
determine its first $t$ digits. The first $t$ digits confine $x$ to a short
interval, and the last $t$ digits confine $x$ to a few residue classes
modulo $g^t$. Write $X_1=g^{L/k}$, so that $x<X_1$, and note that
$X_1\asymp_{g,k}(\#\cP_{g,L})^{2/k}$. Each pair (interval, residue class)
contains at most $X_1g^{-2t}+1$ admissible $x$, roughly, and there are about
$g^t$ pairs. The choice $g^t\asymp X_1^{1/2}$ gives \eqref{eq:CLS}. We are not
aware of any earlier improvement of the exponent $1/k$ in \eqref{eq:CLS},
for any base.

\subsection{Results}
Our main result is that \eqref{eq:CLS} can be improved by a power, in
every base. Let $e(x)=e^{2\pi ix}$ and $F_g(\theta)=\sum_{d=0}^{g-1}e(d\theta)$.
For $j\ge1$ and $0\le n<g^j$ let $\rev_j(n)$ be the integer obtained by
reversing the string of the $j$ digits of $n$, leading zeros included. For
$s\ge1$ and $\ell,a\in\Z$ put
\[
\Psi_\ell(a)=\sum_{m=0}^{g^s-1}e\Bigl(-\frac{\ell\rev_s(m)+am}{g^s}\Bigr),
\qquad
\Sinf(s)=\sum_{0<|\ell|<g^s}\ \max_{a\in\Z}|\Psi_\ell(a)|.
\]
Here $\Psi_\ell(a)$ is the Fourier transform of the function
$m\mapsto e(-\ell\rev_s(m)/g^s)$ on $\Z/g^s\Z$. Trivially
$\Sinf(s)\le2g^{2s}$. Let $\omega(g)$ be the number of distinct prime factors
of $g$.

\begin{theorem}\label{thm:main}
Let $g\ge2$ and $k\ge2$, and let $0<c<1$ be such that
$\Sinf(s)\ll_gg^{(2-c)s}$ for all $s\ge1$. Then the following hold for all
$L\ge3$.
\begin{enumerate}[label=\textup{(\alph*)}]
\item If $g$ is a prime power, then
\[
\#\cQ^k_{g,L}\ \ll_{g,k,c}\ (\#\cP_{g,L})^{\frac1k\left(1-\frac{c}{4-c}\right)}\log\#\cP_{g,L}.
\]
\item If $c\le\frac{\log2}{2\log g}$, then
\[
\#\cQ^k_{g,L}\ \ll_{g,k,c}\ (\#\cP_{g,L})^{\frac1k\left(1-\frac{c}{(k-1)(4-c)+2}\right)}\bigl(\log\#\cP_{g,L}\bigr)^{\omega(g)+1}.
\]
\end{enumerate}
\end{theorem}

If $\Sinf(s)\ll g^{(2-c)s}$ then the same holds for every smaller $c$, so
the condition in (b) is harmless. Part (b) applies to every base; part (a)
is stronger for prime powers. The difference comes from integers $x$ that
share a large factor with a power of $g$; this matters only when $g$ is
not a prime power (see Remark~\ref{rem:whycomposite}).

\begin{theorem}\label{thm:sup}
For every $g\ge2$ there is $c_g>0$ with $\Sinf(s)\ll_gg^{(2-c_g)s}$ for all
$s\ge1$. For the bases in Table~\ref{tab:c} one may take the values of $c_g$
listed there.
\end{theorem}

\begin{table}[ht]
\caption{Admissible values of $c_g$ in Theorem~\ref{thm:sup} (rounded down).}\label{tab:c}
\small
\begin{tabular}{@{}lccccccccccc@{}}
\toprule
$g$ & 2 & 3 & 4 & 5 & 6 & 7 & 8 & 9 & 10 & 12 & 16\\
\midrule
$c_g$ & 0.1263 & 0.1002 & 0.0811 & 0.0709 & 0.0641 & 0.0597 & 0.0543 & 0.0516 & 0.0496 & 0.0458 & 0.0383\\
\bottomrule
\end{tabular}
\end{table}

The existence of some $c_g>0$ is proved analytically in
Section~\ref{subsec:pairing}. The values in Table~\ref{tab:c} come from a
finite computation, described in Sections~\ref{subsec:windows}
and~\ref{subsec:comp}, whose output is certified by a Collatz--Wielandt
bound. It uses $40$-digit arithmetic where this matters and double
precision elsewhere, with a safety margin that dominates the rounding
errors; the accompanying script reproduces it.
Combining the two theorems (and absorbing the logarithms by rounding) gives
the following.

\begin{corollary}\label{cor:main}
For every $g\ge2$ and $k\ge2$ there is $\gamma=\gamma(g,k)<1$ with
$\#\cQ^k_{g,L}\ll_{g,k}(\#\cP_{g,L})^{\gamma/k}$. For the bases below one
may take the following values.
\begin{center}\small\upshape
\begin{tabular}{@{}lcccccccc@{}}
\toprule
prime power $g$ (every $k$) & 2 & 3 & 4 & 5 & 7 & 8 & 9 & 16\\
\midrule
$\gamma$ & 0.9674 & 0.9744 & 0.9794 & 0.9820 & 0.9849 & 0.9863 & 0.9870 & 0.9904\\
\bottomrule
\end{tabular}

\medskip
\begin{tabular}{@{}lccc@{}}
\toprule
other $g$ & 6 & 10 & 12\\
\midrule
$\gamma$, $k=2$ & 0.9893 & 0.9917 & 0.9924\\
$\gamma$, $k=3$ & 0.9936 & 0.9950 & 0.9954\\
\bottomrule
\end{tabular}
\end{center}
In particular, the number of $L$-digit palindromic squares in base $10$
is $\ll(\#\cP_{10,L})^{0.4959}$.
\end{corollary}

For prime powers the table lists $1-c_g/(4-c_g)$, and for the other bases
$1-c_g/((k-1)(4-c_g)+2)$, rounded up. If $g$ is not a prime power, the saving decays
like $1/k$ as $k$ grows.

\subsection{Outline of the method}
We start from the tree of \cite{CLS}: fix $t$, and for each residue class $b$
modulo $g^t$ (a \emph{node}) the admissible $x$ form an arithmetic
progression $x=x_b+g^tz$, $0\le z<Z_b$, of length about $X_1g^{-2t}$. The
argument of \cite{CLS} only uses the first and last $t$ digits. We also use
the next $s$ digits from both ends. Inside a node, the digits of $x^k$ in
positions $t,\dots,t+s-1$ are an affine function of $z$ modulo $g^s$
(a $g$-adic linearisation), and the digits in positions
$L-t-s,\dots,L-t-1$ are, up to a small error, an affine function of $z$
over $\R$ (an archimedean linearisation). The palindrome condition says
that the second block is the reversal of the first. Detecting it with a
Beurling--Selberg majorant and expanding $e(-\ell\rev_s(n)/g^s)$ in
additive characters reduces the count to geometric series in $z$ with
phases $\theta_{b,\ell,a}=(\ell\delta_b+a\beta_b)/g^s$, weighted by
$|\Psi_\ell(a)|$ (Proposition~\ref{prop:red}). Here $\delta_b$ comes from
the size of $x$ and $\beta_b=kb^{k-1}\bmod g^s$ from its residue.

We then bound the weights by $\max_a|\Psi_\ell(a)|$, and recover the
averaging that is needed from the archimedean phases $\ell\delta_b$ alone:
since $\delta_b$ is a monotone function of the first $t$ digits of $x^k$,
and different nodes have essentially different first digits, these phases
are well distributed modulo $1$ across the nodes
(Lemma~\ref{lem:count}). This reduces everything to $\Sinf(s)$.

A power saving for $\Sinf(s)$ comes from an identity between consecutive
factors of the product formula for $\Psi_\ell(a)$, in which the frequency
$a$ cancels (Lemma~\ref{lem:diff}): a large value of $|\Psi_\ell(a)|$ needs
all factors to be large, and the identity makes this impossible for most
$\ell$, whatever $a$ is. Arguments of this kind go back to Col
\cite{Col}; see also \cite{TP}, \cite[Section~2.3]{DMRSS} and \cite{BhS}. We
give a self-contained treatment, and we sharpen the constants considerably
by observing that the identity determines the whole sequence of factor
arguments from the first one (Lemma~\ref{lem:orbit}). This allows us to
bound the supremum over $a$ on windows of consecutive digits, and to sum
over $\ell$ with a transfer matrix (Proposition~\ref{prop:windows}).

If $g$ is not a prime power, there is one more difficulty: an integer $x$ may be
divisible by a high power of one prime factor of $g$, and for such nodes
the reduction loses. We treat these $x$ separately, by a direct count
(Lemma~\ref{lem:large}).

\subsection{Related work}
Uniform bounds for exponential sums involving digit reversal appear in
Col \cite{Col}, Tuxanidy and Panario \cite{TP}, Dartyge, Martin, Rivat,
Shparlinski and Swaenepoel \cite{DMRSS} (in base $2$) and Bhowmik and Suzuki
\cite{BhS} (in general bases, in a preprint). In base $2$, the mechanism
behind Lemma~\ref{lem:diff} appears in \cite[Lemmas~2.5 and~2.6]{DMRSS},
where it is credited to an idea of Col \cite[proof of Lemme~2]{Col}. In
general bases, the identity of Lemma~\ref{lem:diff}, in the form dual to
ours under Lemma~\ref{lem:psi}(iv), appears in the proof of
\cite[Lemma~4.6]{BhS}, a lemma which Bhowmik and Suzuki attribute to
Tuxanidy and Panario \cite{TP}; they
combine it with a pairing of consecutive factors \cite[Lemmas~4.7
and~4.9]{BhS}, as we do in Lemma~\ref{lem:chain}. We do not use any result
of \cite{BhS}: everything we need is proved here. Johnston and Kerr
\cite{JK} (also a preprint) combine the argument of \cite{CLS} with a
hybrid $p$-adic/archimedean van der Corput method to prove that every base
has infinitely many square-free palindromes. They count palindromes
$n\le x$ coprime to $g^3-g$ and divisible by $d^2$ for some $d\in[D,2D]$;
their bounds that go beyond the elementary method of \cite{CLS} require
$D\le x^{2/5}$, while their elementary bound, for $D\asymp x^{1/2}$,
recovers \eqref{eq:CLS} for squares. Square and cube divisors of
palindromes also play a role in the recent preprints of Chourasiya and
Johnston \cite{CJ} and of Tuxanidy \cite{Tux}. We do not use the results
of these preprints.

\subsection{Notation}
We use Vinogradov's notation $\ll$ and $\gg$, and $A\asymp B$ for
$A\ll B\ll A$; implied constants may depend on $g$ and $k$, and on other
parameters only when indicated. We write $\|x\|$ for the distance from $x$
to the nearest integer, $v_p$ for the $p$-adic valuation, and
$0.d_1d_2\cdots d_r$ for $\sum_jd_jg^{-j}$. Throughout,
$X_1=g^{L/k}$; then $x^k$ has length $L$ if and only if
$g^{(L-1)/k}\le x<X_1$, and $X_1\asymp(\#\cP_{g,L})^{2/k}$. Finally
\[
f(\theta)=\frac{|F_g(\theta)|}{g}=\Bigl|\frac{\sin\pi g\theta}{g\sin\pi\theta}\Bigr|\quad(\theta\notin\Z),\qquad f(\theta)=1\quad(\theta\in\Z).
\]
Thus $0\le f\le1$, $f$ is even and $1$-periodic, $f(\theta)<1$ for
$\theta\notin\Z$, and $f$ is continuous.

\section{Preliminaries}

\begin{lemma}\label{lem:psi}
Let $s\ge1$ and $\ell,a\in\Z$.
\begin{enumerate}[label=\textup{(\roman*)}]
\item For $0\le n<g^s$, $e(-\ell\rev_s(n)/g^s)=g^{-s}\sum_{a\bmod g^s}\Psi_\ell(a)e(an/g^s)$.
\item $|\Psi_\ell(a)|=\prod_{i=0}^{s-1}|F_g(\theta_i)|=g^s\prod_{i=0}^{s-1}f(\theta_i)$, where
$\theta_i=\theta_i(\ell,a)=(\ell g^{s-1-i}+ag^i)/g^s$.
\item If $\ell\not\equiv0\pmod{g^s}$ then $\Psi_\ell(0)=0$.
\item $|\Psi_\ell(a)|=|\Psi_a(\ell)|$.
\end{enumerate}
\end{lemma}

\begin{proof}
(i) is Fourier inversion on $\Z/g^s\Z$. (ii) Write $m=\sum_{i<s}m_ig^i$; then
$\rev_s(m)=\sum_im_ig^{s-1-i}$, and $\Psi_\ell(a)=\prod_{i<s}F_g(-\theta_i)$;
note $|F_g(-\theta)|=|F_g(\theta)|$. (iii) Since $\rev_s$ permutes
$[0,g^s)$, $\Psi_\ell(0)=\sum_{n<g^s}e(-\ell n/g^s)=0$. (iv) Substitute
$m=\rev_s(m')$ in the definition; $\rev_s$ is an involution of $[0,g^s)$.
\end{proof}

\begin{lemma}[Beurling--Selberg majorant]\label{lem:selberg}
Let $J\subset\R/\Z$ be an arc of length $|J|<1$ and let $\kappa\ge0$ be an
integer. There is a real trigonometric polynomial
$S(x)=\sum_{|\ell|\le\kappa}\hat S(\ell)e(\ell x)$ with $S\ge\mathbf 1_J$ and
$\hat S(0)=|J|+1/(\kappa+1)$. Moreover $|\hat S(\ell)|\le\hat S(0)$ for all
$\ell$.
\end{lemma}

\begin{proof}
The first statement is Selberg's construction; see \cite{Va} or
\cite[Chapter~1]{Mo}. Since $S\ge0$, $|\hat S(\ell)|\le\int_0^1S=\hat S(0)$.
\end{proof}

\begin{lemma}\label{lem:geom}
For $\theta\in\R$ and an integer $Z\ge1$,
$\bigl|\sum_{z=0}^{Z-1}e(z\theta)\bigr|\le\min\bigl(Z,\frac1{2\|\theta\|}\bigr)$.
\end{lemma}

\begin{proof}
For $\theta\notin\Z$ the sum is $(e(Z\theta)-1)/(e(\theta)-1)$, of modulus at
most $1/|\sin\pi\theta|\le1/(2\|\theta\|)$.
\end{proof}

\begin{lemma}\label{lem:roots}
For every integer $n\ge1$, the group $(\Z/n\Z)^\times$ has at most
$(2k)^{\omega(n)}$ elements $u$ with $u^k=1$.
\end{lemma}

\begin{proof}
By the Chinese remainder theorem it suffices to treat $n=p^a$. For odd $p$,
or $p^a\le4$, the group $(\Z/p^a\Z)^\times$ is cyclic and has at most $k$
such elements; for $a\ge3$, $(\Z/2^a\Z)^\times\cong\Z/2\times\Z/2^{a-2}$ has
at most $2k$ (see \cite[Chapter~4]{IR}).
\end{proof}

\section{The tree and the reduction}\label{sec:reduction}

\subsection{Nodes}
Fix integers $t\ge3$ and $1\le s\le t-2$ with $t+s\le L/2$. A \emph{node}
is an integer $b\in[0,g^t)$ with $g\nmid b^k$. For a node $b$ put
\[
T_b=\rev_t(b^k\bmod g^t),\qquad
I_b=\bigl[(T_bg^{L-t})^{1/k},\ ((T_b+1)g^{L-t})^{1/k}\bigr),
\]
and let $Z_b$ be the number of integers $x\in I_b$ with $x\equiv b\pmod{g^t}$.
If $Z_b\ge1$ let $x_b$ be the least of them, and put
\[
\beta_b=kb^{k-1}\bmod g^s,\qquad\delta_b=kx_b^{k-1}g^{2t+s-L},\qquad
\theta_{b,\ell,a}=\frac{\ell\delta_b+a\beta_b}{g^s}.
\]
Since $g\nmid b^k$, the last digit of $b^k$ is non-zero, so
$g^{t-1}\le T_b<g^t$. Finally put
\[
Z=X_1g^{-2t},\qquad \delta(T)=k(Tg^{L-t})^{(k-1)/k}g^{2t+s-L}\quad(g^{t-1}\le T<g^t).
\]

\begin{lemma}[\cite{CLS}]\label{lem:tree}
If $x^k\in\cP_{g,L}$, then $b=x\bmod g^t$ is a node, $x\in I_b$, and
$x=x_b+g^tz$ for some $0\le z<Z_b$.
\end{lemma}

\begin{proof}
The last digit of a palindrome equals its first digit, which is non-zero;
so $g\nmid x^k$, and hence $g\nmid b^k$. The number formed by the first $t$
digits of $x^k$ is $\lfloor x^k/g^{L-t}\rfloor$, and by the palindrome
property it equals $\rev_t(x^k\bmod g^t)=\rev_t(b^k\bmod g^t)=T_b$. Hence
$T_bg^{L-t}\le x^k<(T_b+1)g^{L-t}$, i.e.\ $x\in I_b$.
\end{proof}

\begin{lemma}\label{lem:sizes}
For every node, $|I_b|g^{-t}\in[Z/k,\ gZ/k]$, so $Z/k-1\le Z_b\le gZ/k+1$;
in particular $Z_b\ge1$ if $Z\ge2k$. If $Z_b\ge1$ then
$0\le\delta_b-\delta(T_b)\le g^{1-t}\delta(T_b)$. Finally
$\frac{k}{2g}\,\frac{g^s}Z\le\delta(T)\le k\,\frac{g^s}Z$ for
$g^{t-1}\le T<g^t$.
\end{lemma}

\begin{proof}
Let $g^{t-1}\le T<g^t$, $y_0=Tg^{L-t}$ and $y_1=y_0+g^{L-t}$. By the mean value
theorem $y_1^{1/k}-y_0^{1/k}=g^{L-t}/(k\xi^{(k-1)/k})$ for some
$\xi\in(y_0,y_1)\subset[g^{L-1},g^L]$, which gives the first claim (take
$T=T_b$). Next, $x_b\in[y_0^{1/k},y_1^{1/k})$ for $T=T_b$, so
$\delta_b/\delta(T_b)=(x_b/y_0^{1/k})^{k-1}\in[1,(y_1/y_0)^{(k-1)/k}]$, and
$y_1/y_0=1+1/T_b\le1+g^{1-t}$. Finally
$\delta(T)\cdot g^{L-t}/(k\xi^{(k-1)/k})=g^{t+s}(y_0/\xi)^{(k-1)/k}\in[g^{t+s}/2,g^{t+s}]$,
because $y_0/\xi\ge y_0/y_1\ge\frac12$; combine with
$g^{L-t}/(k\xi^{(k-1)/k})\in[g^tZ/k,g^{t+1}Z/k]$.
\end{proof}

\subsection{Linearisation}
For a node $b$ with $Z_b\ge1$ put
$\alpha_b=\lfloor x_b^k/g^t\rfloor\bmod g^s$, and for real $z\ge0$
\[
\varphi_b(z)=\frac{(x_b+g^tz)^k}{g^{L-t-s}}-T_bg^s,\qquad
\lambda_b(z)=\varphi_b(0)+\delta_bz .
\]

\begin{lemma}\label{lem:lin}
Let $b$ be a node and $x=x_b+g^tz$ with $0\le z<Z_b$.
\begin{enumerate}[label=\textup{(\alph*)}]
\item The digits of $x^k$ in positions $t,\dots,t+s-1$ form the number
$n_b(z):=(\alpha_b+\beta_bz)\bmod g^s$.
\item The digits of $x^k$ in positions $L-t-s,\dots,L-t-1$ form the number
$\lfloor\varphi_b(z)\rfloor$, and
$0\le\varphi_b(z)-\lambda_b(z)\le\frac{k-1}kg^{s+1-t}\le\frac1g$.
\item If $x^k\in\cP_{g,L}$, then
$\lambda_b(z)-\rev_s(n_b(z))\in[-\frac1g,1)$.
\end{enumerate}
\end{lemma}

\begin{proof}
(a) Write $x^k=x_b^k+kx_b^{k-1}g^tz+g^{2t}N'$ with $N'\in\Z$. Then
$\lfloor x^k/g^t\rfloor=\lfloor x_b^k/g^t\rfloor+kx_b^{k-1}z+g^tN'$; reduce modulo
$g^s$, using $s\le t$ and $x_b\equiv b\pmod{g^t}$.

(b) For $x\in I_b$ we have $\lfloor x^k/g^{L-t}\rfloor=T_b$, so the block in
question is $\lfloor x^k/g^{L-t-s}\rfloor-T_bg^s=\lfloor\varphi_b(z)\rfloor$. As a
function of real $z$, $\varphi_b$ is convex with $\varphi_b'(0)=\delta_b$, so
$\varphi_b\ge\lambda_b$, and by Taylor's theorem
\[
\varphi_b(z)-\lambda_b(z)\le\tfrac12k(k-1)\,x_{\max}^{k-2}(g^tz)^2g^{-(L-t-s)},
\]
where $x_{\max}=\sup I_b$. With $y_0,y_1$ as in the proof of
Lemma~\ref{lem:sizes} (for $T=T_b$) we have $g^tz\le|I_b|\le g^{L-t}/(ky_0^{1-1/k})$
and $x_{\max}\le y_1^{1/k}$, so
\[
x_{\max}^{k-2}(g^tz)^2\le\Bigl(\frac{y_1}{y_0}\Bigr)^{\frac{k-2}k}\frac{g^{2(L-t)}}{k^2y_0}\le\frac{2g^{L-2t+1}}{k^2},
\]
using $y_0\ge g^{L-1}$ and $y_1\le2y_0$. This gives the bound
$\frac{k-1}kg^{s+1-t}$, which is at most $1/g$ because $s\le t-2$.

(c) Since $t+s\le L/2$, the two blocks are disjoint, and the palindrome
property says that the digit in position $L-1-j$ equals the digit in
position $j$. Hence the number formed by the block in (b) is the reversal of
the number formed by the block in (a): $\lfloor\varphi_b(z)\rfloor=\rev_s(n_b(z))$.
So $\rev_s(n_b(z))\in(\varphi_b(z)-1,\varphi_b(z)]$, and the claim follows
from (b).
\end{proof}

\subsection{The reduction}
For a set $\cB$ of nodes with $Z_b\ge1$ for all $b\in\cB$, put
\[
E(\cB)=\sum_{b\in\cB}\ \sum_{0<|\ell|<g^s}\ \sum_{a\not\equiv0\ (g^s)}|\Psi_\ell(a)|\min\Bigl(Z_b,\frac1{2\|\theta_{b,\ell,a}\|}\Bigr).
\]

\begin{proposition}\label{prop:red}
For every such $\cB$,
\[
\#\{x:\ x^k\in\cP_{g,L},\ x\bmod g^t\in\cB\}\ \le\ \frac52g^{-s}\sum_{b\in\cB}Z_b+\frac52g^{-2s}E(\cB).
\]
\end{proposition}

\begin{proof}
Let $J$ be the image of $[-\frac12,1]$ in $\R/g^s\Z$. Applying
Lemma~\ref{lem:selberg} to the arc $g^{-s}J$ of $\R/\Z$ with $\kappa=g^s-1$ and
rescaling, we get $F(y)=\sum_{|\ell|\le\kappa}c_\ell e(\ell y/g^s)$ with
$F\ge\mathbf1_J$ on $\R$, $c_0=\frac{3/2}{g^s}+\frac1{g^s}=\frac5{2g^s}$ and
$|c_\ell|\le c_0$.

Fix $b\in\cB$. By Lemmas~\ref{lem:tree} and~\ref{lem:lin}(c), the number of
solutions $x$ with $x\equiv b\pmod{g^t}$ is at most
\[
\sum_{z<Z_b}F\bigl(\lambda_b(z)-\rev_s(n_b(z))\bigr)
=\sum_{|\ell|\le\kappa}c_\ell\sum_{z<Z_b}e\Bigl(\frac{\ell\lambda_b(z)}{g^s}\Bigr)e\Bigl(-\frac{\ell\rev_s(n_b(z))}{g^s}\Bigr).
\]
Expand the last factor by Lemma~\ref{lem:psi}(i). Since
$e(an_b(z)/g^s)=e(a(\alpha_b+\beta_bz)/g^s)$ and
$\lambda_b(z)=\lambda_b(0)+\delta_bz$, the sum over $z$ becomes
$g^{-s}\sum_a\Psi_\ell(a)\,e(\cdot)\sum_{z<Z_b}e(z\theta_{b,\ell,a})$, with a
unimodular factor $e(\cdot)$ independent of $z$. By Lemma~\ref{lem:geom}
its modulus is at most
$g^{-s}\sum_a|\Psi_\ell(a)|\min(Z_b,1/(2\|\theta_{b,\ell,a}\|))$. For
$\ell=0$, $\Psi_0(a)=g^s\mathbf1_{a\equiv0}$ and the term is $c_0Z_b$. For
$0<|\ell|\le\kappa$ we have $\ell\not\equiv0\pmod{g^s}$, so the term $a\equiv0$
vanishes by Lemma~\ref{lem:psi}(iii). Sum over $b\in\cB$ and use
$|c_\ell|\le c_0$.
\end{proof}

\section{The supremum over the \texorpdfstring{$g$}{g}-adic frequency}\label{sec:count}

We now bound the weights $|\Psi_\ell(a)|$ in $E(\cB)$ by their maximum
over $a$, sum the remaining factor over $a$ (Lemma~\ref{lem:asum}), and
then over the nodes (Lemma~\ref{lem:count}).

\begin{lemma}\label{lem:asum}
Let $y\in\R$, $\beta\in\Z$, $d=\gcd(\beta,g^s)$ and $A\ge1$. Then
\[
\sum_{a\bmod g^s}\min\Bigl(A,\frac1{2\|(y+a\beta)/g^s\|}\Bigr)\ \le\
\min\Bigl(dA,\frac{g^s}{2\|y/d\|}\Bigr)+g^s(2+s\log g).
\]
\end{lemma}

\begin{proof}
As $a$ runs modulo $g^s$, $a\beta$ runs over the multiples of $d$ modulo
$g^s$, each taken $d$ times. Put $Q=g^s/d$ and $y/d=n_0+\epsilon_0$ with
$n_0\in\Z$, $|\epsilon_0|=\|y/d\|$. Then $(y+cd)/g^s=(n_0+c+\epsilon_0)/Q$, and as $c$ runs
modulo $Q$ the residue $j\equiv n_0+c$ runs over $(-Q/2,Q/2]$. For $j=0$ the
distance to $\Z$ is $|\epsilon_0|/Q$; for $j\ne0$ it is at least $(|j|-\frac12)/Q$.
Hence the sum is at most
$d\min(A,Q/(2|\epsilon_0|))+d\sum_{j\ne0}Q/(2|j|-1)\le\min(dA,g^s/(2|\epsilon_0|))+g^s(2+\log Q)$.
\end{proof}

For a non-zero integer $n$ let $m_g(n)=\prod_{p\mid g}p^{v_p(n)}$ be its
\emph{$g$-part}, so that $n/m_g(n)$ is coprime to $g$. For $D\ge1$ let
$\tau_g(D)=\#\{m\le D:\ m\mid g^\infty\}\le(1+\log_2D)^{\omega(g)}$.

\begin{lemma}[nodes with small $g$-part]\label{lem:nodes}
Let $D\ge1$ with $D^k<2^t$, and let $b$ be a node with $m:=m_g(b)\le D$.
\begin{enumerate}[label=\textup{(\roman*)}]
\item $d_b:=\gcd(\beta_b,g^s)$ divides $km^{k-1}$.
\item Every node $b'$ with $T_{b'}=T_b$ has $m_g(b')=m$, and there are at
most $(2k)^{\omega(g)}m^{k-1}$ such nodes.
\end{enumerate}
If $g=p^e$ is a prime power, every node has $m_g(b)<g^{1/k}$.
\end{lemma}

\begin{proof}
Write $b=mu$ with $\gcd(u,g)=1$. (i) $\beta_b\equiv km^{k-1}u^{k-1}\pmod{g^s}$,
so $d_b=\gcd(km^{k-1},g^s)$. (ii) Let $p\mid g$, $e_p=v_p(g)$ and $v=v_p(b)$.
Then $p^{kv}\le m^k<2^t\le p^{e_pt}$, so $kv<e_pt$. If $T_{b'}=T_b$ then
$b'^k\equiv b^k\pmod{p^{e_pt}}$, and the right side has $p$-adic valuation
exactly $kv<e_pt$; hence $v_p(b')=v$. So $m_g(b')=m$, and $b'=mu'$ with
$\gcd(u',g)=1$ and $u'^k\equiv u^k\pmod{g^t/m^k}$ (note $m^k\mid g^t$). By
Lemma~\ref{lem:roots}, $u'$ lies in at most $(2k)^{\omega(g)}$ classes
modulo $g^t/m^k$. Since $b'\bmod g^t$ is determined by $u'\bmod g^t/m$, and
each class modulo $g^t/m^k$ has $m^{k-1}$ lifts modulo $g^t/m$, there are at
most $(2k)^{\omega(g)}m^{k-1}$ such $b'$. Finally, if $g=p^e$ then
$p^e\nmid b^k$ gives $kv_p(b)<e$.
\end{proof}

\begin{lemma}[archimedean count]\label{lem:count}
Let $s\le t-2$, $Z\ge2k$, and $D\ge1$ with $D^k<2^t$. Let $\cB$ be a set of
nodes with $m_g(b)\le D$ for all $b\in\cB$, and for $0<|\ell|<g^s$ put
\[
W_\ell(\cB)=\sum_{b\in\cB}\ \sum_{a\not\equiv0}\min\Bigl(Z_b,\frac1{2\|\theta_{b,\ell,a}\|}\Bigr).
\]
Then, with $\tau=\tau_g(kD^{k-1})$,
\[
W_\ell(\cB)\ll\tau D^{k-1}\Bigl((\log DZ)\,g^{t+s}+g^{2s}+D^{k-1}Z\Bigr)+\tau D^{k-2}\,\frac{Zg^t}{|\ell|}+(s+1)\,g^{t+s}.
\]
\end{lemma}

\begin{proof}
By Lemma~\ref{lem:sizes}, $Z_b\ge1$ for every node. By Lemma~\ref{lem:asum}
with $A=Z_b\le Z_+:=gZ/k+1$,
\[
W_\ell(\cB)\le\sum_{b\in\cB}\min\Bigl(d_bZ_+,\frac{g^s}{2\|\ell\delta_b/d_b\|}\Bigr)+g^{t+s}(2+s\log g).
\]
Split $\cB$ into $\cB_1=\{b:|\ell|\delta_b\le d_b/2\}$ and
$\cB_2=\cB\setminus\cB_1$.

\emph{Small phases.} For $b\in\cB_1$ we have
$\|\ell\delta_b/d_b\|=|\ell|\delta_b/d_b$, and $\delta_b\ge\delta(T_b)\ge
kg^s/(2gZ)$ by Lemma~\ref{lem:sizes}. So the term for $b$ is at most
$gZd_b/(k|\ell|)$. By Lemma~\ref{lem:nodes}(i), and since at most $g^t/m$
residues modulo $g^t$ are divisible by a given $m\mid g^t$,
\[
\sum_{b\in\cB}d_b\le\sum_{\substack{m\le D\\ m\mid g^\infty}}\frac{g^t}m\,km^{k-1}\le k\,\tau_g(D)\,D^{k-2}g^t .
\]
This gives the term $\tau D^{k-2}Zg^t/|\ell|$, since $\tau_g(D)\le\tau$.

\emph{Large phases: reduction to a count over $T$.} Let $b\in\cB_2$. By
Lemma~\ref{lem:sizes} and $s\le t-2$,
$|\ell\delta_b-\ell\delta(T_b)|/d_b\le g^s\cdot\frac{kg^s}Z\cdot\frac{g^{1-t}}{d_b}\le\frac{kg^{s-1}}{d_bZ}=:\varepsilon_1$.
Replace $\delta_b$ by $\delta(T_b)$ in the minimum. If
$\|\ell\delta(T_b)/d_b\|\ge2\varepsilon_1$, the old minimum is at most twice the
new one. Otherwise the new minimum is at least
$\min(d_bZ_+,g^s/(4\varepsilon_1))\ge gd_bZ/(4k)$, while the old one
is at most $d_bZ_+\le gd_bZ/k+d_b$; so again the ratio is at most
$4+4k/(gZ)\le6$. Also
$|\ell|\delta(T_b)\ge|\ell|\delta_b/(1+g^{1-t})>d_b/3$. Group the nodes of
$\cB_2$ by $d=d_b$; by Lemma~\ref{lem:nodes}(i) there are at most $\tau$
values of $d$, all at most $kD^{k-1}$. By Lemma~\ref{lem:nodes}(ii), each
$T$ equals $T_b$ for at most $(2k)^{\omega(g)}D^{k-1}$ nodes $b\in\cB$.
Hence
\[
\sum_{b\in\cB_2}\min\Bigl(d_bZ_+,\frac{g^s}{2\|\ell\delta_b/d_b\|}\Bigr)
\ll D^{k-1}\sum_{d}\ \sum_{T\in\cT_d}\min\Bigl(dZ_+,\frac{g^s}{2\|h_d(T)\|}\Bigr),
\]
where $h_d(T)=\ell\delta(T)/d$ and
$\cT_d=\{T\in[g^{t-1},g^t):|h_d(T)|\ge\frac13\}$, an interval.

\emph{The count.} On $[g^{t-1},g^t)$ the function $h_d$ is monotone with
$|h_d'(T)|=\frac{k-1}k|h_d(T)|/T$, and $|h_d'|$ varies by a factor at most
$g$ there, since $|h_d'|\propto T^{-1/k}$. Let $\mu=\min_{\cT_d}|h_d'|$; then
$\mu\ge\frac{k-1}{3k}g^{-t}$. The range of $h_d$ on $\cT_d$ has length
$R\le g^{t+1}\mu$, and also $R\le|\ell|\max\delta/d\le k|\ell|g^s/(dZ)$. For
every $u>0$, the $T\in\cT_d$ with $\|h_d(T)\|\le u$ lie in at most $R+2$
intervals of length at most $2u/\mu$, each containing at most $2u/\mu+1$
integers, so
\[
n(u):=\#\{T\in\cT_d:\|h_d(T)\|\le u\}\le(R+2)\Bigl(\frac{2u}\mu+1\Bigr)\le2ug^{t+1}+\frac{4u}\mu+R+2\ll ug^{t+1}+R+2.
\]
Split the $T$ dyadically according to $\|h_d(T)\|$, starting at $H/A'$ with
$A'=dZ_+$ and $H=g^s/2$; there are $j_0+2$ ranges with $j_0\ll\log(DZ)$. This
gives
\[
\sum_{T\in\cT_d}\min\Bigl(A',\frac H{\|h_d(T)\|}\Bigr)\le\sum_{j=-1}^{j_0}\frac{A'}{2^j}\,n\Bigl(\frac{2^{j+1}H}{A'}\Bigr)
\ll(j_0+2)\,g^{t+1}H+A'(R+2).
\]
Here $A'R\ll|\ell|g^s\le g^{2s}$ and $A'\le kD^{k-1}Z_+$. Summing over the
at most $\tau$ values of $d$ gives the claim.
\end{proof}

\begin{corollary}\label{cor:E}
Under the hypotheses of Lemma~\ref{lem:count},
\[
E(\cB)\ \ll\ \tau D^{k-1}\Bigl((\log DZ)\,g^{t+s}+g^{2s}+D^{k-1}Z\Bigr)\Sinf(s)
+\tau D^{k-2}(s+1)\,Zg^{t+s}+(s+1)\,g^{t+s}\Sinf(s).
\]
\end{corollary}

\begin{proof}
Clearly $E(\cB)\le\sum_{0<|\ell|<g^s}\max_a|\Psi_\ell(a)|\,W_\ell(\cB)$.
Insert Lemma~\ref{lem:count}. In the term with $1/|\ell|$, use
$\max_a|\Psi_\ell(a)|\le g^s$ and
$\sum_{0<|\ell|<g^s}|\ell|^{-1}\le2(1+s\log g)$.
\end{proof}

\section{The supremum norm of the reversal transform}\label{sec:sup}

\subsection{The identity}
Recall from Lemma~\ref{lem:psi}(ii) that $|\Psi_\ell(a)|=g^s\prod_{i<s}f(\theta_i)$
with $\theta_i=(\ell g^{s-1-i}+ag^i)/g^s$.

\begin{lemma}\label{lem:diff}
For all $\ell,a$ and $0\le i\le s-2$,
$g\theta_i-\theta_{i+1}=(g^2-1)\ell/g^{i+2}$. Consequently
\[
g\|\theta_i\|+\|\theta_{i+1}\|\ \ge\ \Delta_i(\ell):=\Bigl\|\frac{(g^2-1)\ell}{g^{i+2}}\Bigr\|.
\]
\end{lemma}

\begin{proof}
$g\theta_i=\ell/g^i+a/g^{s-1-i}$ and $\theta_{i+1}=\ell/g^{i+2}+a/g^{s-1-i}$.
Subtract.
\end{proof}

The point is that $a$ cancels: a large value of $|\Psi_\ell(a)|$ needs every
$\theta_i$ to be close to an integer, and this is impossible when
$\Delta_i(\ell)$ is not small, whatever $a$ is. In fact the identity
determines the whole sequence $(\theta_i)$ from $\theta_0$:

\begin{lemma}[forced orbit]\label{lem:orbit}
Let $\ell'=(g^2-1)\ell\bmod g^s$ and, for $0\le i\le s-2$,
$D_i=\ell'/g^{i+2}\bmod1=0.\ell'_{i+1}\ell'_i\cdots\ell'_0$, where
$\ell'_j$ are the digits of $\ell'$. Then
$\Delta_i(\ell)=\|D_i\|$ and $\theta_{i+1}\equiv g\theta_i-D_i\pmod1$ for every
$a$. The map $\ell\mapsto\ell'$ is a bijection of $\Z/g^s\Z$.
\end{lemma}

\begin{proof}
By Lemma~\ref{lem:diff}, $\theta_{i+1}=g\theta_i-(g^2-1)\ell/g^{i+2}$, and
$(g^2-1)\ell\equiv\ell'\pmod{g^s}$ with $i+2\le s$. Since
$\gcd(g^2-1,g)=1$, multiplication by $g^2-1$ is a bijection modulo $g^s$.
\end{proof}

\subsection{A power saving in every base}\label{subsec:pairing}
Let $\psi(\epsilon)=\sup\{f(\theta):\|\theta\|\ge\epsilon\}$,
$\phi=-\log\psi\in[0,\infty]$, and
\[
\chi(\Delta)=\inf\{\phi(x)+\phi(y):\ x,y\in[0,\tfrac12],\ gx+y\ge\Delta\}.
\]
The function $\psi$ is non-increasing, and $\phi$ and $\chi$ are
non-decreasing. Since $f$ is continuous and $f<1$ off
$\Z$, we have $\psi(\epsilon)<1$, i.e.\ $\phi(\epsilon)>0$, for
$\epsilon>0$; and $\chi(\Delta)\ge\phi(\Delta/(g+1))>0$ for $\Delta>0$,
because $gx+y\ge\Delta$ forces $\max(x,y)\ge\Delta/(g+1)$.

\begin{lemma}\label{lem:chain}
For every $\ell$,
$\max_a|\Psi_\ell(a)|\le g^s\exp\bigl(-\frac12\sum_{i=0}^{s-2}\chi(\Delta_i(\ell))\bigr)$.
\end{lemma}

\begin{proof}
We have $f(\theta_i)\le\psi(\|\theta_i\|)$, so
$\log(|\Psi_\ell(a)|/g^s)\le-\sum_{i<s}\phi(\|\theta_i\|)$. Each index occurs
in at most two of the pairs $\{i,i+1\}$, $0\le i\le s-2$, so
$\sum_{i<s}\phi(\|\theta_i\|)\ge\frac12\sum_{i\le s-2}\bigl(\phi(\|\theta_i\|)+\phi(\|\theta_{i+1}\|)\bigr)\ge\frac12\sum_{i\le s-2}\chi(\Delta_i(\ell))$
by Lemma~\ref{lem:diff}.
\end{proof}

\begin{proposition}\label{prop:pairing}
For every $g\ge2$ there is $c>0$ with $\Sinf(s)\ll_gg^{(2-c)s}$.
\end{proposition}

\begin{proof}
By Lemma~\ref{lem:orbit} we may sum over $\ell'$ instead of $\ell$; each
non-zero residue is represented by two values of $\ell$ with
$0<|\ell|<g^s$. Let $r=1$ if $g\ge3$ and $r=2$ if $g=2$. By
Lemma~\ref{lem:orbit}, $\Delta_i(\ell)=\|D_i\|$, and
$D_i=0.\ell'_{i+1}\ell'_i\cdots\ell'_0$ lies in the interval $[v,v+g^{-r}]$
with $v=0.\ell'_{i+1}\cdots\ell'_{i+2-r}$. Hence $\Delta_i(\ell)\ge
\Lambda_r(\ell'_{i+1},\dots,\ell'_{i+2-r})$, where
$\Lambda_r(d_1,\dots,d_r)=\max(0,\min(v,1-v-g^{-r}))$ for $v=0.d_1\cdots d_r$.
By Lemma~\ref{lem:chain} and the monotonicity of $\chi$,
\[
\Sinf(s)\le2g^s\sum_{\ell'\bmod g^s}\ \prod_{i=0}^{s-2}\xi(\ell'_{i+1},\dots,\ell'_{i+2-r}),\qquad
\xi(d_1,\dots,d_r)=e^{-\chi(\Lambda_r(d_1,\dots,d_r))/2}.
\]
If $g\ge3$ and $r=1$, the product is over single digits, and the sum is at
most $g\rho^{s-1}$ with $\rho=\sum_{d<g}\xi(d)$. Here $\xi(d)<1$ for $1\le d\le
g-2$, since then $\Lambda_1(d)\ge1/g>0$; so $\rho<g$. If $g=2$ and $r=2$, then
$\Lambda_2(d_1,d_2)=\frac14$ when $d_1\ne d_2$, so the sum is at most
$2\|M^{s-1}\|_\infty$ for the matrix $M=\bigl(\begin{smallmatrix}1&u\\u&1\end{smallmatrix}\bigr)$
with $u=e^{-\chi(1/4)/2}<1$, whose spectral radius is $1+u<2$. In both
cases $\Sinf(s)\ll g^s\rho^s$ with $\rho<g$, which is the claim with
$c=1-\log_g\rho$.
\end{proof}

\begin{example}\label{ex:hand}
For $g=2$, $f(\theta)=|\cos\pi\theta|$ and $\phi(x)=-\log\cos\pi x$. From
$2x+y\ge\frac14$ we get $\max(x,y)\ge\frac1{12}$, so
$\chi(\frac14)\ge-\log\cos(\pi/12)>0.0346$, $u<e^{-0.0173}$,
$1+u<1.9829$, and the proof gives $c_2>0.012$.
\end{example}

\subsection{Windows}\label{subsec:windows}
Lemma~\ref{lem:chain} uses Lemma~\ref{lem:diff} only as an inequality for
each pair of neighbouring factors, and then halves. Lemma~\ref{lem:orbit}
allows us to do much better: we bound the supremum over all real starting
phases on windows of consecutive factors, along the exact recursion.
Transfer matrices with certified Perron--Frobenius bounds are a standard
device for products over digits with a single frequency; the point here is
that the supremum over the second frequency $a$ can be taken window by
window.

Fix a window length $w\ge2$ and weights $\lambda_0,\dots,\lambda_{w-1}\ge0$
with $K=\sum_j\lambda_j>0$. For real $D_0,\dots,D_{w-2}$ put
\[
\cU(D_0,\dots,D_{w-2})=\sup_{\theta\in\R}\ \prod_{j=0}^{w-1}f(\theta_j)^{\lambda_j},\qquad
\theta_0=\theta,\quad\theta_{j+1}=g\theta_j-D_j .
\]

\begin{lemma}\label{lem:window}
If $w\le s$, then for every $\ell$, with $D_i$ as in Lemma~\ref{lem:orbit},
\[
\max_a|\Psi_\ell(a)|\le g^s\prod_{i=0}^{s-w}\cU(D_i,\dots,D_{i+w-2})^{1/K}.
\]
\end{lemma}

\begin{proof}
Fix $a$. By Lemma~\ref{lem:orbit} and the periodicity of $f$, for each
$0\le i\le s-w$ the numbers $\theta_i,\dots,\theta_{i+w-1}$ give the value of
a product as in the definition of $\cU(D_i,\dots,D_{i+w-2})$, so
$\sum_j\lambda_j\log f(\theta_{i+j})\le\log\cU(D_i,\dots,D_{i+w-2})$. Sum over
$i$. Each index $m\in[0,s)$ receives total weight
$\mathrm{cov}_m\le K$, and $\log f\le0$, so
$K\sum_m\log f(\theta_m)\le\sum_m\mathrm{cov}_m\log f(\theta_m)\le\sum_i\log\cU(\cdots)$.
Finally use Lemma~\ref{lem:psi}(ii).
\end{proof}

\begin{proposition}\label{prop:windows}
Fix $w$ and $\lambda$ as above and an integer $r\ge1$, and put $\nu=w+r-2$;
assume that $\nu\ge2$. For
$\sigma=(\sigma_0,\dots,\sigma_{\nu-1})\in\{0,\dots,g-1\}^\nu$ let
$v(\sigma)=\sum_{j=0}^{r-1}\sigma_{r-1-j}g^{-j-1}$ and
\[
\widehat U(\sigma)=\sup\Bigl\{\cU(D_0,\dots,D_{w-2}):\
D_0\in\bigl[v(\sigma),v(\sigma)+g^{-r}\bigr],\ D_m=\frac{\sigma_{r-1+m}+D_{m-1}}g\Bigr\},
\]
where $1\le m\le w-2$. Let $\overline U(\sigma)\ge\widehat U(\sigma)$ for all
$\sigma$, and let $M$ be the non-negative matrix indexed by
$\{0,\dots,g-1\}^{\nu-1}$ with
\[
M\bigl[(\sigma_0,\dots,\sigma_{\nu-2}),(\sigma_1,\dots,\sigma_{\nu-1})\bigr]=\overline U(\sigma)^{1/K}
\]
and all other entries $0$. If $\mathbf x$ is a vector with positive entries
and $\rho_{\mathbf x}=\max_i(M\mathbf x)_i/\mathbf x_i$, then
$\Sinf(s)\ll g^s\rho_{\mathbf x}^s$, where the implied constant depends on
$g,w,r,\lambda,\overline U$ and $\mathbf x$. That is, Theorem~\ref{thm:sup} holds with
$c_g=1-\log_g\rho_{\mathbf x}$.
\end{proposition}

\begin{proof}
For $s<w$ the claim follows from the trivial bound $\Sinf(s)\le2g^{2s}$, so
let $s\ge w$. For the window starting at $i$ put
$\sigma^{(i)}=(\ell'_{i+2-r},\dots,\ell'_{i+w-1})$, with $\ell'_j=0$ for
$j<0$. By Lemma~\ref{lem:orbit}, the number $D_i$ lies in the interval
$[v(\sigma^{(i)}),v(\sigma^{(i)})+g^{-r}]$, and
$D_{i+m}=(\ell'_{i+m+1}+D_{i+m-1})/g$; hence
$\cU(D_i,\dots,D_{i+w-2})\le\overline U(\sigma^{(i)})$. As in the proof of
Proposition~\ref{prop:pairing}, Lemmas~\ref{lem:window} and~\ref{lem:orbit}
give
\[
\Sinf(s)\le2g^s\sum_{\ell'\bmod g^s}\ \prod_{i=0}^{s-w}\overline U(\sigma^{(i)})^{1/K}.
\]
Consecutive $\sigma^{(i)}$ overlap in $\nu-1$ digits, so each product is
the weight of a path with $n=s-w+1$ edges in the graph of $M$. Summing also
over the (fictitious) digits $\ell'_j$, $j<0$, and over $\ell'_0$ when $r=1$
(it then lies in no window), we see that the sum over $\ell'$ is at most
$g\,\mathbf 1^{\mathsf T}M^n\mathbf 1\le g^\nu\max_i(M^n\mathbf 1)_i$, where
$\mathbf 1$ is the all-ones vector. Finally, $M\mathbf x\le\rho_{\mathbf x}\mathbf x$
entrywise and $M\ge0$ give $M^n\mathbf x\le\rho_{\mathbf x}^n\mathbf x$, hence
$M^n\mathbf 1\le(\max_i\mathbf x_i/\min_i\mathbf x_i)\,\rho_{\mathbf x}^n\mathbf 1$.
\end{proof}

The weights $\overline U(\sigma)^{1/K}$ are positive, since $\widehat U(\sigma)>0$
($f$ vanishes only at finitely many points modulo $1$), and the graph of
$M$ is a de~Bruijn graph, which is strongly connected. So $M$ is
irreducible, and by the Collatz--Wielandt formula the infimum of
$\rho_{\mathbf x}$ over all positive $\mathbf x$ is the spectral radius of
$M$ \cite[Chapter~8]{HJ}. In practice, power iteration produces a vector
$\mathbf x$ for which $\rho_{\mathbf x}$ is very close to the spectral radius.

\subsection{Computing the constants}\label{subsec:comp}
We now describe how the values in Table~\ref{tab:c} were obtained. A Python
script that carries out the computation accompanies this paper (see
Section~\ref{subsec:checks}).

\emph{Upper bounds for $\widehat U$.} Parametrise an orbit by its last point.
Put $C_0=f^{\lambda_0}$ and
\[
C_{m+1}(y)=f(y)^{\lambda_{m+1}}\max_{0\le q<g}\ \sup_{D_m}\ C_m\Bigl(\frac{y+D_m+q}g\Bigr)\qquad(0\le m\le w-2),
\]
where $D_m$ runs over the interval $[v_m,v_m+g^{-r-m}]$, with $v_0=v(\sigma)$
and $v_m=(\sigma_{r-1+m}+v_{m-1})/g$. Then $\widehat U(\sigma)\le\sup_yC_{w-1}(y)$.
Indeed, the numbers $(y+D_m+q)/g$, $0\le q<g$, are the preimages of $y$
modulo $1$ under $\theta\mapsto g\theta-D_m$, so by induction $C_m(y)$ is at
least $\prod_{j\le m}f(\theta_j)^{\lambda_j}$ for every orbit with
$\theta_m\equiv y\pmod1$; treating the $D_m$ as independent in their
intervals only enlarges the bound.

Since the argument of $C_m$ is contracted by $g$, a fixed grid suffices. We
use the closed cells $[n/N,(n+1)/N]$, $0\le n<N$, and the following upper
bound for $f$ on each cell. It is $1$ on the two cells that contain an
integer. On the other cells, $|\sin\pi\theta|$ is minimal at an endpoint,
and $|\sin\pi g\theta|$ is at most $1$, and is maximal at an endpoint if the
cell contains no $\theta$ with $g\theta\in\frac12+\Z$. These bounds are
propagated through the recursion: the bound for $C_{m+1}$ on a cell is the
bound for $f^{\lambda_{m+1}}$ on that cell times the maximum of the bounds
for $C_m$ over blocks of consecutive cells that contain the preimages of the
cell, for all $D_m$ in its interval. These blocks are computed in exact
integer arithmetic, and no Lipschitz constants are needed.

\emph{Rounding.} The cell bounds for $f$ are evaluated in $40$-digit
decimal arithmetic, multiplied by $1+10^{-12}$, and rounded upward to
double precision. All later steps are carried out in double precision,
where maxima are exact; we assume, as holds for standard implementations,
that each power, product, sum or quotient has relative error at most
$2^{-52}$ (any bound below $10^{-14}$ would do). The script also checks
that every computed $\overline U(\sigma)$ exceeds $10^{-280}$, so underflow
plays no role. Consider one step of the recursion on one cell. If the cell
bound for $f$ equals $1$, the step consists of maxima and a multiplication
by $1$, which are exact. Otherwise the cell bound for $f$ exceeds the exact
one by a factor of about $1+10^{-12}$; since $\lambda_{m+1}\ge1$ (all our
weights are positive integers), this factor covers the rounding errors of
the power and of the product. By induction, the computed bounds for $C_m$
are at least the exact ones, so the computed values $\overline U(\sigma)$
are upper bounds for $\widehat U(\sigma)$. In the last stage, each ratio
$(M\mathbf x)_i/\mathbf x_i$ involves $3g$ roundings ($g$ powers, $g$
products, $g-1$ sums and one quotient) and the rounding of the exponent
$1/K$, whose relative effect on an entry $e$ is at most
$|\log e|\,2^{-53}<10^{-13}/2$, because $e>10^{-140}$. The total relative
effect is below $10^{-13}$, and we multiply the final maximum by
$1+10^{-12}$.

\emph{Parameters.} We use triangular weights $\lambda_j=\min(j+1,w-j)$,
which reduce the edge effects of the windows, $N=8192$, and the vector
$\mathbf x$ produced by power iteration. Table~\ref{tab:rho} lists the
parameters and the resulting upper bounds $\rho_{\rm CW}$ for
$\rho_{\mathbf x}$; then
$c_g=1-\log_g\rho_{\rm CW}$, rounded down, is the value in
Table~\ref{tab:c}. The cost grows like $g^{w+r-1}N$, which limits $w$ for
large $g$.

\begin{table}[ht]
\caption{Window parameters, certified bounds $\rho_{\rm CW}$ (rounded up)
and the resulting values of $c_g$ (rounded down). All rows use triangular
weights and $N=8192$.}\label{tab:rho}
\small
\begin{tabular}{@{}rrrll@{}}
\toprule
$g$ & $w$ & $r$ & $\rho_{\rm CW}$ & $c_g$\\
\midrule
2 & 14 & 5 & 1.832340 & 0.1263\\
3 & 10 & 3 & 2.687053 & 0.1002\\
4 & 7 & 3 & 3.574206 & 0.0811\\
5 & 6 & 3 & 4.460721 & 0.0709\\
6 & 6 & 2 & 5.348980 & 0.0641\\
7 & 6 & 2 & 6.231057 & 0.0597\\
8 & 5 & 2 & 7.145309 & 0.0543\\
9 & 5 & 2 & 8.034354 & 0.0516\\
10 & 6 & 1 & 8.919963 & 0.0496\\
12 & 5 & 2 & 10.706920 & 0.0458\\
16 & 4 & 2 & 14.387430 & 0.0383\\
\bottomrule
\end{tabular}
\end{table}

\begin{proof}[Proof of Theorem~\ref{thm:sup}]
The existence of $c_g>0$ for every $g$ is Proposition~\ref{prop:pairing}.
The values in Table~\ref{tab:c} follow from Proposition~\ref{prop:windows}
and the computation just described.
\end{proof}

\begin{remark}\label{rem:sharp}
For $g=2$ and $r=5$, uniform weights give $c_2\ge0.114$, $0.117$, $0.119$
for $w=10,12,14$, and triangular weights give $0.121$, $0.124$, $0.126$.
Numerically, $\Sinf(s+1)/\Sinf(s)$ settles at about $2^{1.866}$
($s\le16$), which corresponds to $c_2\approx0.134$; for $g=3$ and $g=4$ the
corresponding values are about $0.108$ and $0.090$, and for $g=10$ the
ratio at $s=4$ corresponds to $c_{10}\approx0.064$. If these values
reflect the true growth of $\Sinf$, then no bound for $\Sinf$ can improve
$\gamma$ much further for $g=2$: with $c_2=0.134$,
Theorem~\ref{thm:main}(a) would give $\gamma\approx0.965$, against our
$0.9674$. The remaining gap comes partly from the window length and partly
from the discretisation. The pairing bound of
Proposition~\ref{prop:pairing}, with more digits $r$ and the same
discretisation, gives only $c_2\approx0.06$ and $c_{10}\approx0.03$.
\end{remark}

\begin{remark}
The power saving for $\Sinf$ uses the special structure of digit reversal
only through the product formula of Lemma~\ref{lem:psi}(ii) and the
identity of Lemma~\ref{lem:diff}. For a general permutation $\pi$ of
$[0,g^s)$, let $\Psi^\pi_\ell(a)$ be defined as $\Psi_\ell(a)$ with $\pi$
in place of $\rev_s$. Then $\max_a|\Psi^\pi_\ell(a)|$ can equal $g^s$ for
every $\ell$, so that the analogue of $\Sinf(s)$ is of order $g^{2s}$; for
instance, $\pi(m)=hm\bmod g^s$ with $\gcd(h,g)=1$ gives
$|\Psi^\pi_\ell(-h\ell)|=g^s$.
\end{remark}

\section{Proof of \texorpdfstring{Theorem~\ref{thm:main}}{Theorem 1.1}}\label{sec:proof}

\begin{remark}\label{rem:whycomposite}
For a prime power $g$, Lemma~\ref{lem:nodes} applies to every node with
$D=g^{1/k}$, a constant. If $g$ is not a prime power, a node can be divisible by a
high power of one prime factor of $g$ and still satisfy $g\nmid b^k$ (in
base $10$, for example $b=2^v$). For such nodes both $d_b$ and the
multiplicity of $b\mapsto T_b$ are about $m_g(b)^{k-1}$, and the
corresponding progressions resonate together. We therefore treat the
integers $x$ with large $g$-part separately, by a direct count.
\end{remark}

\begin{lemma}[large $g$-part]\label{lem:large}
Let $P$ be the largest prime factor of $g$, and let $D\ge1$ with
$PD\le X_1^{\log2/(2\log g)}$. Then
\[
\#\{x:\ x^k\in\cP_{g,L},\ m_g(x)>D\}\ \ll\ \tau_g(PD)\,X_1^{1/2}D^{-1}.
\]
\end{lemma}

\begin{proof}
Multiplying in the prime factors of $m_g(x)$ one at a time, we find a
divisor $m'$ of $m_g(x)$ with $D<m'\le PD$. Fix such an $m'$, and choose $t'$
with $X_1^{1/2}\le g^{t'}<gX_1^{1/2}$. Then
$p^{v_p(m')}\le PD\le2^{t'}\le p^{v_p(g)t'}$ for each $p\mid g$, so
$m'\mid g^{t'}$. Apply Lemma~\ref{lem:tree} with $t'$ in place of $t$ (this
part of Lemma~\ref{lem:tree} does not use $s$): the $x$ in question with
$m'\mid x$ lie in nodes $b$ of level $t'$ with $m'\mid b$, of which there are at
most $g^{t'}/m'$. By the first assertion of Lemma~\ref{lem:sizes} (which
does not involve $s$, and where now $X_1g^{-2t'}\le1$), each such node
contains at most $g/k+1$ of them. So there are $\ll X_1^{1/2}/m'\le
X_1^{1/2}/D$ such $x$, and there are at most $\tau_g(PD)$ choices of $m'$.
\end{proof}

\begin{proof}[Proof of Theorem~\ref{thm:main}]
Recall $X_1\asymp(\#\cP_{g,L})^{2/k}$, and note
$\frac2k\cdot\frac{2-c}{4-c}=\frac1k\bigl(1-\frac c{4-c}\bigr)$. We may assume
that $L$ is large.

\emph{(a)} Let $g=p^e$, $D=g^{1/k}$, and let $\cB$ be the set of all nodes;
by Lemma~\ref{lem:nodes} every node has $m_g(b)<D$, and $\tau\ll1$. Choose
$t$ with $g^{(4-c)t}\asymp X_1$ and put $s=t-2$. Then
$Z=X_1g^{-2t}\asymp g^{(2-c)t}$; for large $L$ we have $2k\le Z$,
$D^k=g<2^t$, $D^{k-1}Z\le g^{2s}$, and $t+s<2t\le L/2$ because
$t=L/(k(4-c))+O(1)$ and $k(4-c)>6$. By Lemma~\ref{lem:tree} every solution
$x$ satisfies $x\bmod g^t\in\cB$. By Proposition~\ref{prop:red},
Lemma~\ref{lem:sizes} and Corollary~\ref{cor:E}, using $g^{t-s}=g^2$ and
$s\asymp\log Z\asymp\log X_1$,
\[
\#\cQ^k_{g,L}\ll Z+1+(\log X_1)\bigl(\Sinf(s)+Z\bigr)\ll(\log X_1)\,X_1^{\frac{2-c}{4-c}} ,
\]
since $\Sinf(s)\ll g^{(2-c)s}\le g^{(2-c)t}\asymp Z\asymp X_1^{(2-c)/(4-c)}$.

\emph{(b)} Let
\[
\eta=\frac{c}{2\bigl((k-1)(4-c)+2\bigr)},\qquad D=X_1^\eta .
\]
Split the $x$ with $x^k\in\cQ^k_{g,L}$ according to whether $m_g(x)>D$.
Those with $m_g(x)>D$ number $\ll(\log X_1)^{\omega(g)}X_1^{1/2-\eta}$ by
Lemma~\ref{lem:large}; its condition holds because
$\eta<c/4\le\log2/(8\log g)$.

For the others, choose $t$ with $g^{(4-c)t}\asymp X_1/D$, put $s=t-2$, and
let $\cB$ be the set of nodes $b$ with $m_g(b)\le D$. Then
$t=(1-\eta)L/(k(4-c))+O(1)$, so for large $L$ we have $t\ge3$ and
$t+s<2t\le L/2$. Moreover $D^k<2^t$ for large $L$, because
\[
k\eta\le\frac c{4-c}\le\frac{\log2}{2(4-c)\log g}<(1-\eta)\frac{\log2}{(4-c)\log g}
\quad\text{and}\quad 2^t\gg X_1^{(1-\eta)\log2/((4-c)\log g)} .
\]
If $m_g(x)\le D$ then $p^{v_p(x)}\le D<2^t\le p^{v_p(g)t}$ for each
$p\mid g$, so $m_g(x\bmod g^t)=m_g(x)$, and by Lemma~\ref{lem:tree},
$x\bmod g^t\in\cB$. Now $Z=X_1g^{-2t}\asymp X_1^{\frac{2-c}{4-c}}D^{\frac2{4-c}}$
and $Dg^{(2-c)t}\asymp X_1^{\frac{2-c}{4-c}}D^{\frac2{4-c}}$ as well; in
particular $Z\ge2k$ for large $L$, so Lemma~\ref{lem:count} applies. We
shall also use $D^{k-1}Z\le g^{2s}$, which holds for large $L$ because
$\eta\bigl((k-1)(4-c)+4\bigr)<c$. By Proposition~\ref{prop:red},
Lemma~\ref{lem:sizes} and Corollary~\ref{cor:E}, with
$\tau\ll(\log X_1)^{\omega(g)}$ and $g^{t-s}=g^2$,
\begin{align*}
\#\{x:\ m_g(x)\le D\}&\ll Z+\tau(\log X_1)\Bigl(D^{k-1}\Sinf(s)+D^{2k-2}Zg^{-2s}\Sinf(s)+D^{k-2}Z\Bigr)\\
&\ll(\log X_1)^{\omega(g)+1}D^{k-2}\bigl(D\,g^{(2-c)t}+Z\bigr),
\end{align*}
using $D^{k-1}Z\le g^{2s}$ in the middle term. The right side is
$\ll(\log X_1)^{\omega(g)+1}D^{k-2+\frac2{4-c}}X_1^{\frac{2-c}{4-c}}$. With
$D=X_1^\eta$ the exponent of $X_1$ is
\[
\eta\Bigl(k-2+\frac2{4-c}\Bigr)+\frac{2-c}{4-c}=\frac12-\eta,
\]
by the choice of $\eta$. Finally
$2(\frac12-\eta)=1-\frac c{(k-1)(4-c)+2}$.
\end{proof}

\begin{proof}[Proof of Corollary~\ref{cor:main}]
The first statement follows from Theorems~\ref{thm:main} and~\ref{thm:sup},
replacing $c$ by $\min(c_g,\frac{\log2}{2\log g})$ in part (b). For the
table, apply Theorem~\ref{thm:main}(a) for prime powers and (b) for $g=6,10,12$,
with $c=c_g$ from Table~\ref{tab:c}; these satisfy
$c_g\le\log2/(2\log g)$. The table entries are the resulting exponents
rounded up in the fourth decimal, which absorbs the logarithmic factors.
\end{proof}

\section{Further remarks}\label{sec:remarks}

\subsection{Numerical checks}\label{subsec:checks}
The script \texttt{window\_constants.py}, which accompanies this paper,
implements Section~\ref{subsec:comp}: for each base it prints
$\rho_{\rm CW}$ and $c_g$, and it reproduces Table~\ref{tab:rho} in about
an hour on one processor core. In a second mode it computes, for
small $s$, the values of $\max_a|\Psi_\ell(a)|$ by the fast Fourier
transform (in floating point), and compares them for every $\ell$ with the bound of
Lemma~\ref{lem:window}, with $\overline U$ in place of $\cU$. For
$(g,s,w,r)=(2,12,8,5)$, $(3,7,5,3)$, $(4,7,5,3)$, $(6,5,4,2)$,
$(10,4,3,1)$, $(10,4,3,2)$ and $(16,3,3,1)$, with $N=4096$, the bound held
for every $\ell$; the largest ratio was $0.99995$, for $g=16$. Finally, we
checked
the identity of Lemma~\ref{lem:diff} exhaustively for small $(g,s)$,
Lemma~\ref{lem:orbit} in exact rational arithmetic, and the structural
statements of Lemma~\ref{lem:nodes} for all nodes for
$(g,k,t)=(10,2,4)$, $(10,2,5)$, $(10,3,4)$, $(6,2,5)$, $(6,3,5)$ and
$(12,2,4)$; these checks are not needed for the proofs.

\subsection{Limits of the method}
The argument bounds the weights $|\Psi_\ell(a)|$ by their maximum over $a$
and uses averaging only over the archimedean phases. By
Remark~\ref{rem:sharp}, this is unlikely to give much more than
$\gamma\approx0.965$ in base $2$. Going further would require cancellation in the sum over $a$,
that is, information on the joint distribution of the archimedean phases
$\ell\delta_b$ and the $g$-adic phases $a\beta_b$ across the nodes. For
bases that are not prime powers, the treatment of integers with a large $g$-part
(Lemma~\ref{lem:large}) is also wasteful and could probably be improved.

\section*{Acknowledgements}
The arguments and computations in this paper were developed with the
assistance of Claude, an AI model made by Anthropic. Claude did not have
access to the version of
\cite{DMRSS} published in the Journal of the London Mathematical Society;
references to specific results of \cite{DMRSS} are to its arXiv version,
arXiv:2309.11380v1.

\begin{thebibliography}{99}
\bibitem{BHS} W.~D. Banks, D.~N. Hart and M.~Sakata, \emph{Almost all palindromes are composite}, Math.\ Res.\ Lett.\ \textbf{11} (2004), no.~5--6, 853--868.
\bibitem{BhS} G.~Bhowmik and Y.~Suzuki, \emph{On Telhcirid's theorem on arithmetic progressions}, preprint, arXiv:2406.13334v1 (2024).
\bibitem{CJ} S.~Chourasiya and D.~R. Johnston, \emph{Power-free palindromes and reversed primes}, preprint, arXiv:2503.21136v5 (2025).
\bibitem{CLB} J.~Cilleruelo, F.~Luca and L.~Baxter, \emph{Every positive integer is a sum of three palindromes}, Math.\ Comp.\ \textbf{87} (2018), no.~314, 3023--3055.
\bibitem{CLS} J.~Cilleruelo, F.~Luca and I.~E. Shparlinski, \emph{Power values of palindromes}, J.\ Comb.\ Number Theory \textbf{1} (2009), no.~2, 101--107.
\bibitem{Col} S.~Col, \emph{Palindromes dans les progressions arithm\'etiques}, Acta Arith.\ \textbf{137} (2009), no.~1, 1--41.
\bibitem{DMRSS} C.~Dartyge, B.~Martin, J.~Rivat, I.~E. Shparlinski and C.~Swaenepoel, \emph{Reversible primes}, J.\ Lond.\ Math.\ Soc.\ (2) \textbf{109} (2024), no.~3, Paper No.\ e12883; arXiv:2309.11380.
\bibitem{HJ} R.~A. Horn and C.~R. Johnson, \emph{Matrix Analysis}, 2nd ed., Cambridge University Press, Cambridge, 2013.
\bibitem{IR} K.~Ireland and M.~Rosen, \emph{A Classical Introduction to Modern Number Theory}, 2nd ed., Graduate Texts in Mathematics \textbf{84}, Springer, New York, 1990.
\bibitem{JK} D.~R. Johnston and B.~Kerr, \emph{The infinitude of square-free palindromes}, preprint, arXiv:2601.07097v2 (2026).
\bibitem{Mo} H.~L. Montgomery, \emph{Ten Lectures on the Interface between Analytic Number Theory and Harmonic Analysis}, CBMS Regional Conference Series in Mathematics \textbf{84}, American Mathematical Society, Providence, RI, 1994.
\bibitem{Tux} A.~Tuxanidy, \emph{Distribution of square-free palindromes}, preprint, arXiv:2603.02549v2 (2026).
\bibitem{TP} A.~Tuxanidy and D.~Panario, \emph{Infinitude of palindromic almost-prime numbers}, Int.\ Math.\ Res.\ Not.\ IMRN \textbf{2024}, no.~18, 12466--12503.
\bibitem{Va} J.~D. Vaaler, \emph{Some extremal functions in Fourier analysis}, Bull.\ Amer.\ Math.\ Soc.\ (N.S.) \textbf{12} (1985), no.~2, 183--216.
\end{thebibliography}

\end{document}
